\begin{table*}[t]
		\caption{两种不同计算方法（即基于占用\cite{hu2022st}和基于稀疏）在自车真实未来轨迹上计算的碰撞率比较。}
		\begin{center}
			\tiny
			% \renewcommand\arraystretch{1.5}
			\vspace{-10pt}
	%		\resizebox{1\textwidth}{!}
			{
				\begin{tabular}{p{3cm}|p{1.5cm}<{\centering}|p{0.7cm}<{\centering}p{0.7cm}<{\centering}p{0.7cm}<{\centering}p{0.7cm}<{\centering}}
					\toprule
					\multirow{2}{*}{指标}& \multirow{2}{*}{分辨率}& \multicolumn{4}{c}{碰撞率(\%)} \\
					&& 1s& 2s& 3s& 平均\\
					\midrule
					Occupancy-Based\cite{hu2022st}&0.5m& 10.39	&10.33&	10.25& 10.32\\
					Sparse-Based(本文)&-&0.58&	0.58&	0.58& 0.58\\
					\bottomrule
				\end{tabular}% }
			}
		\end{center}
		% \vspace{2pt}
		% \vspace{-20pt}
		\label{tab:gt-collision}
	\end{table*}

	\begin{table*}[t]
		\caption{通过所提出的基于稀疏的碰撞度量对不同方法规划结果的碰撞率比较。$\ddag$：使用官方检查点复现的结果。}
		\begin{center}
			\tiny
			% \renewcommand\arraystretch{1.5}
			\vspace{-10pt}
			%		\resizebox{1\textwidth}{!}
			{
				\begin{tabular}{c|c|cccc|cccc}
					\toprule
					\multirow{2}{*}{指标}& \multirow{2}{*}{碰撞优化器}& \multicolumn{4}{c|}{L2(m)}& \multicolumn{4}{c}{碰撞率(\%)} \\
					&&1s& 2s& 3s& 平均& 1s& 2s& 3s& 平均\\
					\midrule
					\multirow{2}{*}{UniAD$\ddag$\cite{hu2023planning}}&\Checkmark&0.44&0.69&1.03 &0.72&1.07&	1.35&	1.62&1.35\\
					&\XSolidBrush&0.31&0.57&0.93&0.60&0.52&	0.73&	1.19&0.81\\
					\midrule
					{SparseAD-B}&-&0.15&0.31&0.56&0.29&	0.00& 0.04&	0.14&0.06\\
					\bottomrule
				\end{tabular}% }
			}
		\end{center}
		% \vspace{2pt}
		% \vspace{-20pt}
		\label{tab:pred-collision}
	\end{table*}